% Options for packages loaded elsewhere
\PassOptionsToPackage{unicode}{hyperref}
\PassOptionsToPackage{hyphens}{url}
\documentclass[
]{article}
\usepackage{xcolor}
\usepackage[margin=1in]{geometry}
\usepackage{amsmath,amssymb}
\setcounter{secnumdepth}{-\maxdimen} % remove section numbering
\usepackage{iftex}
\ifPDFTeX
  \usepackage[T1]{fontenc}
  \usepackage[utf8]{inputenc}
  \usepackage{textcomp} % provide euro and other symbols
\else % if luatex or xetex
  \usepackage{unicode-math} % this also loads fontspec
  \defaultfontfeatures{Scale=MatchLowercase}
  \defaultfontfeatures[\rmfamily]{Ligatures=TeX,Scale=1}
\fi
\usepackage{lmodern}
\ifPDFTeX\else
  % xetex/luatex font selection
\fi
% Use upquote if available, for straight quotes in verbatim environments
\IfFileExists{upquote.sty}{\usepackage{upquote}}{}
\IfFileExists{microtype.sty}{% use microtype if available
  \usepackage[]{microtype}
  \UseMicrotypeSet[protrusion]{basicmath} % disable protrusion for tt fonts
}{}
\makeatletter
\@ifundefined{KOMAClassName}{% if non-KOMA class
  \IfFileExists{parskip.sty}{%
    \usepackage{parskip}
  }{% else
    \setlength{\parindent}{0pt}
    \setlength{\parskip}{6pt plus 2pt minus 1pt}}
}{% if KOMA class
  \KOMAoptions{parskip=half}}
\makeatother
\usepackage{color}
\usepackage{fancyvrb}
\newcommand{\VerbBar}{|}
\newcommand{\VERB}{\Verb[commandchars=\\\{\}]}
\DefineVerbatimEnvironment{Highlighting}{Verbatim}{commandchars=\\\{\}}
% Add ',fontsize=\small' for more characters per line
\usepackage{framed}
\definecolor{shadecolor}{RGB}{248,248,248}
\newenvironment{Shaded}{\begin{snugshade}}{\end{snugshade}}
\newcommand{\AlertTok}[1]{\textcolor[rgb]{0.94,0.16,0.16}{#1}}
\newcommand{\AnnotationTok}[1]{\textcolor[rgb]{0.56,0.35,0.01}{\textbf{\textit{#1}}}}
\newcommand{\AttributeTok}[1]{\textcolor[rgb]{0.13,0.29,0.53}{#1}}
\newcommand{\BaseNTok}[1]{\textcolor[rgb]{0.00,0.00,0.81}{#1}}
\newcommand{\BuiltInTok}[1]{#1}
\newcommand{\CharTok}[1]{\textcolor[rgb]{0.31,0.60,0.02}{#1}}
\newcommand{\CommentTok}[1]{\textcolor[rgb]{0.56,0.35,0.01}{\textit{#1}}}
\newcommand{\CommentVarTok}[1]{\textcolor[rgb]{0.56,0.35,0.01}{\textbf{\textit{#1}}}}
\newcommand{\ConstantTok}[1]{\textcolor[rgb]{0.56,0.35,0.01}{#1}}
\newcommand{\ControlFlowTok}[1]{\textcolor[rgb]{0.13,0.29,0.53}{\textbf{#1}}}
\newcommand{\DataTypeTok}[1]{\textcolor[rgb]{0.13,0.29,0.53}{#1}}
\newcommand{\DecValTok}[1]{\textcolor[rgb]{0.00,0.00,0.81}{#1}}
\newcommand{\DocumentationTok}[1]{\textcolor[rgb]{0.56,0.35,0.01}{\textbf{\textit{#1}}}}
\newcommand{\ErrorTok}[1]{\textcolor[rgb]{0.64,0.00,0.00}{\textbf{#1}}}
\newcommand{\ExtensionTok}[1]{#1}
\newcommand{\FloatTok}[1]{\textcolor[rgb]{0.00,0.00,0.81}{#1}}
\newcommand{\FunctionTok}[1]{\textcolor[rgb]{0.13,0.29,0.53}{\textbf{#1}}}
\newcommand{\ImportTok}[1]{#1}
\newcommand{\InformationTok}[1]{\textcolor[rgb]{0.56,0.35,0.01}{\textbf{\textit{#1}}}}
\newcommand{\KeywordTok}[1]{\textcolor[rgb]{0.13,0.29,0.53}{\textbf{#1}}}
\newcommand{\NormalTok}[1]{#1}
\newcommand{\OperatorTok}[1]{\textcolor[rgb]{0.81,0.36,0.00}{\textbf{#1}}}
\newcommand{\OtherTok}[1]{\textcolor[rgb]{0.56,0.35,0.01}{#1}}
\newcommand{\PreprocessorTok}[1]{\textcolor[rgb]{0.56,0.35,0.01}{\textit{#1}}}
\newcommand{\RegionMarkerTok}[1]{#1}
\newcommand{\SpecialCharTok}[1]{\textcolor[rgb]{0.81,0.36,0.00}{\textbf{#1}}}
\newcommand{\SpecialStringTok}[1]{\textcolor[rgb]{0.31,0.60,0.02}{#1}}
\newcommand{\StringTok}[1]{\textcolor[rgb]{0.31,0.60,0.02}{#1}}
\newcommand{\VariableTok}[1]{\textcolor[rgb]{0.00,0.00,0.00}{#1}}
\newcommand{\VerbatimStringTok}[1]{\textcolor[rgb]{0.31,0.60,0.02}{#1}}
\newcommand{\WarningTok}[1]{\textcolor[rgb]{0.56,0.35,0.01}{\textbf{\textit{#1}}}}
\usepackage{graphicx}
\makeatletter
\newsavebox\pandoc@box
\newcommand*\pandocbounded[1]{% scales image to fit in text height/width
  \sbox\pandoc@box{#1}%
  \Gscale@div\@tempa{\textheight}{\dimexpr\ht\pandoc@box+\dp\pandoc@box\relax}%
  \Gscale@div\@tempb{\linewidth}{\wd\pandoc@box}%
  \ifdim\@tempb\p@<\@tempa\p@\let\@tempa\@tempb\fi% select the smaller of both
  \ifdim\@tempa\p@<\p@\scalebox{\@tempa}{\usebox\pandoc@box}%
  \else\usebox{\pandoc@box}%
  \fi%
}
% Set default figure placement to htbp
\def\fps@figure{htbp}
\makeatother
\setlength{\emergencystretch}{3em} % prevent overfull lines
\providecommand{\tightlist}{%
  \setlength{\itemsep}{0pt}\setlength{\parskip}{0pt}}
\usepackage{booktabs}
\usepackage{longtable}
\usepackage{array}
\usepackage{multirow}
\usepackage{wrapfig}
\usepackage{float}
\usepackage{colortbl}
\usepackage{pdflscape}
\usepackage{tabu}
\usepackage{threeparttable}
\usepackage{threeparttablex}
\usepackage[normalem]{ulem}
\usepackage{makecell}
\usepackage{xcolor}
\usepackage{bookmark}
\IfFileExists{xurl.sty}{\usepackage{xurl}}{} % add URL line breaks if available
\urlstyle{same}
\hypersetup{
  pdftitle={Week\_5\_Practical},
  pdfauthor={Owen Hughes},
  hidelinks,
  pdfcreator={LaTeX via pandoc}}

\title{Week\_5\_Practical}
\author{Owen Hughes}
\date{2025-10-28}

\begin{document}
\maketitle

\textbf{Alright so today I want to make a few maps showing the amount of
hotels and airbnbs in various london boroughs}

\begin{Shaded}
\begin{Highlighting}[]
\FunctionTok{print}\NormalTok{(}\StringTok{"load libraries"}\NormalTok{)}
\end{Highlighting}
\end{Shaded}

\begin{verbatim}
## [1] "load libraries"
\end{verbatim}

\begin{Shaded}
\begin{Highlighting}[]
\FunctionTok{library}\NormalTok{(sf)}
\end{Highlighting}
\end{Shaded}

\begin{verbatim}
## Linking to GEOS 3.13.1, GDAL 3.11.0, PROJ 9.6.0; sf_use_s2() is TRUE
\end{verbatim}

\begin{Shaded}
\begin{Highlighting}[]
\FunctionTok{library}\NormalTok{(here)}
\end{Highlighting}
\end{Shaded}

\begin{verbatim}
## here() starts at C:/Users/Owen Hughes/Desktop/CASA/UCL_Work/GIS_Term1/GIS_Git_Projects/Week_5/GIS
\end{verbatim}

\begin{Shaded}
\begin{Highlighting}[]
\FunctionTok{library}\NormalTok{(janitor)}
\end{Highlighting}
\end{Shaded}

\begin{verbatim}
## 
## Attaching package: 'janitor'
\end{verbatim}

\begin{verbatim}
## The following objects are masked from 'package:stats':
## 
##     chisq.test, fisher.test
\end{verbatim}

\begin{Shaded}
\begin{Highlighting}[]
\FunctionTok{library}\NormalTok{(tidyverse)}
\end{Highlighting}
\end{Shaded}

\begin{verbatim}
## -- Attaching core tidyverse packages ------------------------ tidyverse 2.0.0 --
## v dplyr     1.1.4     v readr     2.1.5
## v forcats   1.0.1     v stringr   1.5.2
## v ggplot2   4.0.0     v tibble    3.3.0
## v lubridate 1.9.4     v tidyr     1.3.1
## v purrr     1.1.0
\end{verbatim}

\begin{verbatim}
## -- Conflicts ------------------------------------------ tidyverse_conflicts() --
## x dplyr::filter() masks stats::filter()
## x dplyr::lag()    masks stats::lag()
## i Use the conflicted package (<http://conflicted.r-lib.org/>) to force all conflicts to become errors
\end{verbatim}

\begin{Shaded}
\begin{Highlighting}[]
\FunctionTok{library}\NormalTok{(tmap)}
\FunctionTok{library}\NormalTok{(ggplot2)}
\end{Highlighting}
\end{Shaded}

\textbf{Right, the specific goal is to make one map showing the hotels
in each London borough with some context, and then another map of the
same size and legend showing airbnbs and its own context}

Okay for the hotel and airbnb maps I will have the boroughs as a shape
file, hotels as a shape file, and airbnbs as a csv. I want to put all of
these in the same crs so that they can be joined later.

To attain the correct information, I will have to filter hotels from the
larger Ldn data, and filter entire home and apartments available for 365
days a year for the airbnb data.

\begin{Shaded}
\begin{Highlighting}[]
\FunctionTok{print}\NormalTok{(}\StringTok{"Time to Load data"}\NormalTok{)}
\end{Highlighting}
\end{Shaded}

\begin{verbatim}
## [1] "Time to Load data"
\end{verbatim}

\begin{Shaded}
\begin{Highlighting}[]
\CommentTok{\# for the hotel data }
\NormalTok{OSM }\OtherTok{\textless{}{-}} \FunctionTok{st\_read}\NormalTok{(}\FunctionTok{here}\NormalTok{(}\StringTok{"Week\_5/Data/Ldn\_data/gis\_osm\_pois\_free\_1.shp"}\NormalTok{))}\SpecialCharTok{\%\textgreater{}\%}
  \FunctionTok{st\_transform}\NormalTok{(., }\DecValTok{27700}\NormalTok{) }\SpecialCharTok{\%\textgreater{}\%}
  \CommentTok{\#select hotels only}
  \FunctionTok{filter}\NormalTok{(fclass }\SpecialCharTok{==} \StringTok{\textquotesingle{}hotel\textquotesingle{}}\NormalTok{)}
\end{Highlighting}
\end{Shaded}

\begin{verbatim}
## Reading layer `gis_osm_pois_free_1' from data source 
##   `C:\Users\Owen Hughes\Desktop\CASA\UCL_Work\GIS_Term1\GIS_Git_Projects\Week_5\GIS\Week_5\Data\Ldn_data\gis_osm_pois_free_1.shp' 
##   using driver `ESRI Shapefile'
## Simple feature collection with 97061 features and 4 fields
## Geometry type: POINT
## Dimension:     XY
## Bounding box:  xmin: -0.5120832 ymin: 51.28679 xmax: 0.3094684 ymax: 51.6902
## Geodetic CRS:  WGS 84
\end{verbatim}

\begin{Shaded}
\begin{Highlighting}[]
\FunctionTok{print}\NormalTok{(}\StringTok{"Time to Load data"}\NormalTok{)}
\end{Highlighting}
\end{Shaded}

\begin{verbatim}
## [1] "Time to Load data"
\end{verbatim}

\begin{Shaded}
\begin{Highlighting}[]
\CommentTok{\# for the airbnb data}
\NormalTok{Airbnb }\OtherTok{\textless{}{-}} \FunctionTok{read\_csv}\NormalTok{(}\FunctionTok{here}\NormalTok{(}\StringTok{"Week\_5/Data/listings.csv"}\NormalTok{)) }\SpecialCharTok{\%\textgreater{}\%}
  \CommentTok{\# longitude is considered x value here, latitude is y}
  \FunctionTok{st\_as\_sf}\NormalTok{(., }\AttributeTok{coords =} \FunctionTok{c}\NormalTok{(}\StringTok{"longitude"}\NormalTok{, }\StringTok{"latitude"}\NormalTok{), }
                   \AttributeTok{crs =} \DecValTok{4326}\NormalTok{) }\SpecialCharTok{\%\textgreater{}\%}
    \FunctionTok{st\_transform}\NormalTok{(., }\DecValTok{27700}\NormalTok{)}\SpecialCharTok{\%\textgreater{}\%}
    \CommentTok{\#select entire places that are available all year}
    \FunctionTok{filter}\NormalTok{(room\_type }\SpecialCharTok{==} \StringTok{\textquotesingle{}Entire home/apt\textquotesingle{}} \SpecialCharTok{\&}\NormalTok{ availability\_365 }\SpecialCharTok{==}\StringTok{\textquotesingle{}365\textquotesingle{}}\NormalTok{)}
\end{Highlighting}
\end{Shaded}

\begin{verbatim}
## Rows: 96651 Columns: 18
## -- Column specification --------------------------------------------------------
## Delimiter: ","
## chr   (4): name, host_name, neighbourhood, room_type
## dbl  (11): id, host_id, latitude, longitude, price, minimum_nights, number_o...
## lgl   (2): neighbourhood_group, license
## date  (1): last_review
## 
## i Use `spec()` to retrieve the full column specification for this data.
## i Specify the column types or set `show_col_types = FALSE` to quiet this message.
\end{verbatim}

\begin{Shaded}
\begin{Highlighting}[]
\FunctionTok{print}\NormalTok{(}\StringTok{"Time to Load data"}\NormalTok{)}
\end{Highlighting}
\end{Shaded}

\begin{verbatim}
## [1] "Time to Load data"
\end{verbatim}

\begin{Shaded}
\begin{Highlighting}[]
\CommentTok{\# For the London Boroughs}
\NormalTok{Londonborough }\OtherTok{\textless{}{-}} \FunctionTok{st\_read}\NormalTok{(}\FunctionTok{here}\NormalTok{(}\StringTok{"Week\_5/Data/statistical{-}gis{-}boundaries{-}london/ESRI/London\_Borough\_Excluding\_MHW.shp"}\NormalTok{))}\SpecialCharTok{\%\textgreater{}\%}
  \FunctionTok{st\_transform}\NormalTok{(., }\DecValTok{27700}\NormalTok{)}
\end{Highlighting}
\end{Shaded}

\begin{verbatim}
## Reading layer `London_Borough_Excluding_MHW' from data source 
##   `C:\Users\Owen Hughes\Desktop\CASA\UCL_Work\GIS_Term1\GIS_Git_Projects\Week_5\GIS\Week_5\Data\statistical-gis-boundaries-london\ESRI\London_Borough_Excluding_MHW.shp' 
##   using driver `ESRI Shapefile'
## Simple feature collection with 33 features and 7 fields
## Geometry type: MULTIPOLYGON
## Dimension:     XY
## Bounding box:  xmin: 503568.2 ymin: 155850.8 xmax: 561957.5 ymax: 200933.9
## Projected CRS: OSGB36 / British National Grid
\end{verbatim}

For now I am going to skip the download for the inset map of the UK
because the data download is not functioning correctly

\textbf{Now it is time to start working on the join of the data sets}

this is how it's done with st\_join

\begin{Shaded}
\begin{Highlighting}[]
\FunctionTok{print}\NormalTok{(}\StringTok{"Perform a left join with st\_join"}\NormalTok{)}
\end{Highlighting}
\end{Shaded}

\begin{verbatim}
## [1] "Perform a left join with st_join"
\end{verbatim}

\begin{Shaded}
\begin{Highlighting}[]
\NormalTok{Hotels }\OtherTok{\textless{}{-}}  \FunctionTok{st\_join}\NormalTok{(Londonborough, OSM)}

\NormalTok{Airbnbs }\OtherTok{\textless{}{-}}  \FunctionTok{st\_join}\NormalTok{(Londonborough, Airbnb)}

\FunctionTok{head}\NormalTok{(Hotels)}
\end{Highlighting}
\end{Shaded}

\begin{verbatim}
## Simple feature collection with 6 features and 11 fields
## Geometry type: MULTIPOLYGON
## Dimension:     XY
## Bounding box:  xmin: 516362.6 ymin: 155850.8 xmax: 539657.9 ymax: 172367
## Projected CRS: OSGB36 / British National Grid
##                     NAME  GSS_CODE HECTARES NONLD_AREA ONS_INNER SUB_2009
## 1   Kingston upon Thames E09000021 3726.117          0         F     <NA>
## 1.1 Kingston upon Thames E09000021 3726.117          0         F     <NA>
## 1.2 Kingston upon Thames E09000021 3726.117          0         F     <NA>
## 1.3 Kingston upon Thames E09000021 3726.117          0         F     <NA>
## 1.4 Kingston upon Thames E09000021 3726.117          0         F     <NA>
## 2                Croydon E09000008 8649.441          0         F     <NA>
##     SUB_2006      osm_id code fclass                             name
## 1       <NA>   387544035 2401  hotel             Bosco Hotel & Lounge
## 1.1     <NA>  1729750835 2401  hotel                       Travelodge
## 1.2     <NA>  4002880204 2401  hotel                      Premier Inn
## 1.3     <NA>  5955546285 2401  hotel The Bull and Bush Hotel Kingston
## 1.4     <NA> 12167570871 2401  hotel                       Travelodge
## 2       <NA>  5308320604 2401  hotel    easyHotel Croydon Town Centre
##                           geometry
## 1   MULTIPOLYGON (((516401.6 16...
## 1.1 MULTIPOLYGON (((516401.6 16...
## 1.2 MULTIPOLYGON (((516401.6 16...
## 1.3 MULTIPOLYGON (((516401.6 16...
## 1.4 MULTIPOLYGON (((516401.6 16...
## 2   MULTIPOLYGON (((535009.2 15...
\end{verbatim}

\begin{Shaded}
\begin{Highlighting}[]
\FunctionTok{view}\NormalTok{(Hotels)}
\FunctionTok{view}\NormalTok{(Airbnbs)}
\end{Highlighting}
\end{Shaded}

Right we can see that each hotel and airbnb in each borough becomes a
new row of the data set.

\textbf{In order to arrange the data so that each row is a single
borough with all of its hotels or airbnbs, we can use the group\_by and
summarize functions}

\begin{Shaded}
\begin{Highlighting}[]
\FunctionTok{print}\NormalTok{(}\StringTok{"Use group by and summarise"}\NormalTok{)}
\end{Highlighting}
\end{Shaded}

\begin{verbatim}
## [1] "Use group by and summarise"
\end{verbatim}

\begin{Shaded}
\begin{Highlighting}[]
\NormalTok{Hotels\_sum }\OtherTok{\textless{}{-}}\NormalTok{ Hotels }\SpecialCharTok{\%\textgreater{}\%}
  \CommentTok{\# we need to list the columns we want to keep in the summarise}
  \FunctionTok{group\_by}\NormalTok{(., GSS\_CODE, NAME)}\SpecialCharTok{\%\textgreater{}\%}
  \CommentTok{\# for each group count the number of rows and store in a column called accomodation count.}
  \FunctionTok{summarise}\NormalTok{(}\StringTok{\textasciigrave{}}\AttributeTok{Accomodation count}\StringTok{\textasciigrave{}} \OtherTok{=} \FunctionTok{n}\NormalTok{())}
\end{Highlighting}
\end{Shaded}

\begin{verbatim}
## `summarise()` has grouped output by 'GSS_CODE'. You can override using the
## `.groups` argument.
\end{verbatim}

\begin{Shaded}
\begin{Highlighting}[]
\NormalTok{Airbnb\_sum }\OtherTok{\textless{}{-}}\NormalTok{ Airbnbs }\SpecialCharTok{\%\textgreater{}\%}
  \FunctionTok{group\_by}\NormalTok{(., GSS\_CODE, NAME)}\SpecialCharTok{\%\textgreater{}\%}
  \FunctionTok{summarise}\NormalTok{(}\StringTok{\textasciigrave{}}\AttributeTok{Accomodation count}\StringTok{\textasciigrave{}} \OtherTok{=} \FunctionTok{n}\NormalTok{())}
\end{Highlighting}
\end{Shaded}

\begin{verbatim}
## `summarise()` has grouped output by 'GSS_CODE'. You can override using the
## `.groups` argument.
\end{verbatim}

\begin{Shaded}
\begin{Highlighting}[]
\FunctionTok{view}\NormalTok{(Hotels\_sum)}
\FunctionTok{view}\NormalTok{(Airbnb\_sum)}

\FunctionTok{print}\NormalTok{(}\StringTok{"The data has been arranges correctly now"}\NormalTok{)}
\end{Highlighting}
\end{Shaded}

\begin{verbatim}
## [1] "The data has been arranges correctly now"
\end{verbatim}

\textbf{The st\_join retains the geometry of the input data set so that
even when there is no hotel or airbnb in a borough, the geometry will
still appear}

\begin{Shaded}
\begin{Highlighting}[]
\FunctionTok{print}\NormalTok{(}\StringTok{"Check and area with no hotels / airbnbs"}\NormalTok{)}
\end{Highlighting}
\end{Shaded}

\begin{verbatim}
## [1] "Check and area with no hotels / airbnbs"
\end{verbatim}

\begin{Shaded}
\begin{Highlighting}[]
\NormalTok{Hotels }\SpecialCharTok{\%\textgreater{}\%}
  \FunctionTok{filter}\NormalTok{(NAME}\SpecialCharTok{==}\StringTok{"Barking and Dagenham"}\NormalTok{)}
\end{Highlighting}
\end{Shaded}

\begin{verbatim}
## Simple feature collection with 1 feature and 11 fields
## Geometry type: MULTIPOLYGON
## Dimension:     XY
## Bounding box:  xmin: 543417.3 ymin: 181434.1 xmax: 551944.6 ymax: 191137.3
## Projected CRS: OSGB36 / British National Grid
##                    NAME  GSS_CODE HECTARES NONLD_AREA ONS_INNER SUB_2009
## 32 Barking and Dagenham E09000002 3779.934     169.15         F     <NA>
##    SUB_2006     osm_id code fclass               name
## 32     <NA> 9090311683 2401  hotel Barking Park Hotel
##                          geometry
## 32 MULTIPOLYGON (((543905.4 18...
\end{verbatim}

\begin{Shaded}
\begin{Highlighting}[]
\NormalTok{Hotels\_sum }\SpecialCharTok{\%\textgreater{}\%}
  \FunctionTok{filter}\NormalTok{(NAME}\SpecialCharTok{==}\StringTok{"Barking and Dagenham"}\NormalTok{)}
\end{Highlighting}
\end{Shaded}

\begin{verbatim}
## Simple feature collection with 1 feature and 3 fields
## Geometry type: MULTIPOLYGON
## Dimension:     XY
## Bounding box:  xmin: 543417.3 ymin: 181434.1 xmax: 551944.6 ymax: 191137.3
## Projected CRS: OSGB36 / British National Grid
## # A tibble: 1 x 4
## # Groups:   GSS_CODE [1]
##   GSS_CODE  NAME                 `Accomodation count`                   geometry
## * <chr>     <chr>                               <int>         <MULTIPOLYGON [m]>
## 1 E09000002 Barking and Dagenham                    1 (((543905.4 183199.1, 543~
\end{verbatim}

\textbf{The problem is, Barking and Dagenham returns a value of 1 when
there is in fact, no hotel in the area. We can use a topological
relationship instead of an st\_join to remedy this}

This way, a count of points inside each polygon is created rather than a
join of the spatial layers

\begin{Shaded}
\begin{Highlighting}[]
\FunctionTok{print}\NormalTok{(}\StringTok{"Topological realtionship with hotels"}\NormalTok{)}
\end{Highlighting}
\end{Shaded}

\begin{verbatim}
## [1] "Topological realtionship with hotels"
\end{verbatim}

\begin{Shaded}
\begin{Highlighting}[]
\NormalTok{Hotels\_example }\OtherTok{\textless{}{-}}\FunctionTok{st\_contains}\NormalTok{(Londonborough, OSM)}

\NormalTok{Hotels\_example}
\end{Highlighting}
\end{Shaded}

\begin{verbatim}
## Sparse geometry binary predicate list of length 33, where the predicate
## was `contains'
## first 10 elements:
##  1: 35, 105, 157, 221, 321
##  2: 213, 223
##  3: 1, 42, 214, 258, 267, 316
##  4: 7, 82, 115, 246, 252, 274
##  5: 39, 239, 250, 254, 291, 320
##  6: 110, 188, 215, 216, 217, 253, 319
##  7: 65, 158, 185, 225, 226, 227, 228, 255
##  8: 145, 162
##  9: 38, 125, 138, 151, 192, 194, 237, 322, 340
##  10: 93, 108, 220, 270, 288
\end{verbatim}

\begin{Shaded}
\begin{Highlighting}[]
\FunctionTok{print}\NormalTok{(}\StringTok{"Find the length of of the list of hotels (features contained in each borough) to get the total"}\NormalTok{)}
\end{Highlighting}
\end{Shaded}

\begin{verbatim}
## [1] "Find the length of of the list of hotels (features contained in each borough) to get the total"
\end{verbatim}

\begin{Shaded}
\begin{Highlighting}[]
\NormalTok{Accomodation\_contained }\OtherTok{\textless{}{-}}\NormalTok{ Londonborough}\SpecialCharTok{\%\textgreater{}\%}
  \FunctionTok{mutate}\NormalTok{(}\AttributeTok{hotels\_n =} \FunctionTok{lengths}\NormalTok{(}\FunctionTok{st\_contains}\NormalTok{(., OSM)))}\SpecialCharTok{\%\textgreater{}\%}
  \FunctionTok{mutate}\NormalTok{(}\AttributeTok{airbnbs\_n =} \FunctionTok{lengths}\NormalTok{(}\FunctionTok{st\_contains}\NormalTok{(., Airbnbs)))}
\end{Highlighting}
\end{Shaded}

\textbf{Now I can check to see that Bexley returns a value of 0 no 1}

\begin{Shaded}
\begin{Highlighting}[]
\FunctionTok{print}\NormalTok{(}\StringTok{"Check Bexley"}\NormalTok{)}
\end{Highlighting}
\end{Shaded}

\begin{verbatim}
## [1] "Check Bexley"
\end{verbatim}

\begin{Shaded}
\begin{Highlighting}[]
\NormalTok{Accomodation\_contained }\SpecialCharTok{\%\textgreater{}\%}
  \FunctionTok{filter}\NormalTok{(NAME}\SpecialCharTok{==}\StringTok{"Bexley"}\NormalTok{)}
\end{Highlighting}
\end{Shaded}

\begin{verbatim}
## Simple feature collection with 1 feature and 9 fields
## Geometry type: MULTIPOLYGON
## Dimension:     XY
## Bounding box:  xmin: 544314.9 ymin: 169908.2 xmax: 554089.2 ymax: 181521.8
## Projected CRS: OSGB36 / British National Grid
##     NAME  GSS_CODE HECTARES NONLD_AREA ONS_INNER SUB_2009 SUB_2006
## 1 Bexley E09000004 6428.649    370.619         F     <NA>     <NA>
##                         geometry hotels_n airbnbs_n
## 1 MULTIPOLYGON (((547226.2 18...        0        39
\end{verbatim}

\textbf{Alright, now that the joins have been completed successfully, I
can start thinking about mapping the data}

Tmap will be used in this practical over ggplot as it is designed
specifically for spatial data

\begin{Shaded}
\begin{Highlighting}[]
\FunctionTok{print}\NormalTok{(}\StringTok{"Start the Tmap"}\NormalTok{)}
\end{Highlighting}
\end{Shaded}

\begin{verbatim}
## [1] "Start the Tmap"
\end{verbatim}

\begin{Shaded}
\begin{Highlighting}[]
\CommentTok{\# standard map {-}{-}\textgreater{} as opposed to interactive}
\FunctionTok{tmap\_mode}\NormalTok{(}\StringTok{"plot"}\NormalTok{)}
\end{Highlighting}
\end{Shaded}

\begin{verbatim}
## i tmap modes "plot" - "view"
## i toggle with `tmap::ttm()`
\end{verbatim}

\begin{Shaded}
\begin{Highlighting}[]
\NormalTok{tm1 }\OtherTok{\textless{}{-}} \FunctionTok{tm\_shape}\NormalTok{(Accomodation\_contained) }\SpecialCharTok{+} 
  \FunctionTok{tm\_polygons}\NormalTok{(}\StringTok{"hotels\_n"}\NormalTok{,}
              \AttributeTok{col =} \StringTok{"\#7f1734"}\NormalTok{, }\AttributeTok{lwd=}\FloatTok{0.5}\NormalTok{, }\AttributeTok{lty=}\StringTok{"dashed"}\NormalTok{,)}

\NormalTok{tm1}
\end{Highlighting}
\end{Shaded}

\pandocbounded{\includegraphics[keepaspectratio]{Week_5_Practical_files/figure-latex/unnamed-chunk-11-1.pdf}}

This would be the map with no spatial data

\begin{Shaded}
\begin{Highlighting}[]
\FunctionTok{print}\NormalTok{(}\StringTok{"plain map"}\NormalTok{)}
\end{Highlighting}
\end{Shaded}

\begin{verbatim}
## [1] "plain map"
\end{verbatim}

\begin{Shaded}
\begin{Highlighting}[]
\NormalTok{tm\_no\_map\_layer }\OtherTok{\textless{}{-}} \FunctionTok{tm\_shape}\NormalTok{(Accomodation\_contained) }\SpecialCharTok{+} 
  \CommentTok{\# there is no tm\_polygons() if we just want the map with no spatial data}
  \FunctionTok{tm\_borders}\NormalTok{(}\AttributeTok{col =} \StringTok{"\#7f1734"}\NormalTok{)}

\NormalTok{tm\_no\_map\_layer}
\end{Highlighting}
\end{Shaded}

\pandocbounded{\includegraphics[keepaspectratio]{Week_5_Practical_files/figure-latex/unnamed-chunk-12-1.pdf}}

In order to add more information to the map, one simply builds off og
the tm1 variable already set

I am going to add a \emph{Violin Plot, Color Scheme, Separation Method
(jenks)}

\begin{Shaded}
\begin{Highlighting}[]
\FunctionTok{print}\NormalTok{(}\StringTok{"Make some additions"}\NormalTok{)}
\end{Highlighting}
\end{Shaded}

\begin{verbatim}
## [1] "Make some additions"
\end{verbatim}

\begin{Shaded}
\begin{Highlighting}[]
\FunctionTok{tmap\_mode}\NormalTok{(}\StringTok{"plot"}\NormalTok{)}
\end{Highlighting}
\end{Shaded}

\begin{verbatim}
## i tmap modes "plot" - "view"
\end{verbatim}

\begin{Shaded}
\begin{Highlighting}[]
\NormalTok{tm1 }\OtherTok{\textless{}{-}} \FunctionTok{tm\_shape}\NormalTok{(Accomodation\_contained) }\SpecialCharTok{+} 
  \FunctionTok{tm\_polygons}\NormalTok{(}\AttributeTok{fill =}\StringTok{"hotels\_n"}\NormalTok{,}
              \AttributeTok{col =} \StringTok{"black"}\NormalTok{, }
              \AttributeTok{lwd =}\FloatTok{0.5}\NormalTok{,}
              \AttributeTok{lty=}\StringTok{"dashed"}\NormalTok{,}
              \AttributeTok{fill.chart =} \FunctionTok{tm\_chart\_violin}\NormalTok{(),}
              \CommentTok{\# above this was the same as before}
              \AttributeTok{fill.scale =} \FunctionTok{tm\_scale\_intervals}\NormalTok{(}
                \AttributeTok{values=}\StringTok{"brewer.bu\_pu"}\NormalTok{,}
                \AttributeTok{n=}\DecValTok{5}\NormalTok{,}
                \AttributeTok{style=}\StringTok{"jenks"}\NormalTok{))}

\NormalTok{tm1}
\end{Highlighting}
\end{Shaded}

\pandocbounded{\includegraphics[keepaspectratio]{Week_5_Practical_files/figure-latex/unnamed-chunk-13-1.pdf}}
What if I wanted to see all the possible options for color schemes?

\begin{Shaded}
\begin{Highlighting}[]
\CommentTok{\#install.packages("colorblindcheck")}
\CommentTok{\# install.packages("shiny")}
\CommentTok{\#install.packages(c("shinyjs", "kableExtra"))}
\FunctionTok{library}\NormalTok{(colorblindcheck)}
\NormalTok{cols4all}\SpecialCharTok{::}\FunctionTok{c4a\_gui}\NormalTok{()}
\end{Highlighting}
\end{Shaded}

\textbf{Right! Now that we have used jenks with the hotel data, we need
to use the same intervals for the airbnb data so that the two maps can
be comparable}

This means I am going to make custom breaks in the data --\textgreater{}
I will derive jenks data for airbnbs and then apply it to hotels

\begin{Shaded}
\begin{Highlighting}[]
\FunctionTok{print}\NormalTok{(}\StringTok{"Sumamry of data"}\NormalTok{)}
\end{Highlighting}
\end{Shaded}

\begin{verbatim}
## [1] "Sumamry of data"
\end{verbatim}

\begin{Shaded}
\begin{Highlighting}[]
\NormalTok{stats }\OtherTok{\textless{}{-}}\NormalTok{ Accomodation\_contained }\SpecialCharTok{\%\textgreater{}\%}
  \FunctionTok{st\_drop\_geometry}\NormalTok{() }\SpecialCharTok{\%\textgreater{}\%}
\NormalTok{  dplyr}\SpecialCharTok{::}\FunctionTok{select}\NormalTok{(hotels\_n, airbnbs\_n) }\SpecialCharTok{\%\textgreater{}\%}  
  \FunctionTok{summarise}\NormalTok{(}\FunctionTok{across}\NormalTok{(}\FunctionTok{everything}\NormalTok{(), }\FunctionTok{list}\NormalTok{(}
    \AttributeTok{min =}\NormalTok{ min,}
    \AttributeTok{max =}\NormalTok{ max,}
    \AttributeTok{mean =}\NormalTok{ mean,}
    \AttributeTok{median =}\NormalTok{ median,}
    \AttributeTok{sd =}\NormalTok{ sd}
\NormalTok{  )))}
\end{Highlighting}
\end{Shaded}

Airbnbs have the largest range so I will use the airbnbs to derive my
data breaks

\begin{Shaded}
\begin{Highlighting}[]
\FunctionTok{print}\NormalTok{(}\StringTok{"derive airbnb jenks"}\NormalTok{)}
\end{Highlighting}
\end{Shaded}

\begin{verbatim}
## [1] "derive airbnb jenks"
\end{verbatim}

\begin{Shaded}
\begin{Highlighting}[]
\FunctionTok{library}\NormalTok{(classInt)}

\NormalTok{breaks }\OtherTok{\textless{}{-}}\NormalTok{ Accomodation\_contained}\SpecialCharTok{\%\textgreater{}\%}
  \FunctionTok{st\_drop\_geometry}\NormalTok{()}\SpecialCharTok{\%\textgreater{}\%}
  \CommentTok{\#need a numeric vector not a dataframe or tibble}
  \FunctionTok{pull}\NormalTok{(airbnbs\_n) }\SpecialCharTok{\%\textgreater{}\%}             
  \FunctionTok{classIntervals}\NormalTok{(., }\AttributeTok{n =} \DecValTok{5}\NormalTok{, }\AttributeTok{style =} \StringTok{"jenks"}\NormalTok{)}
  
\NormalTok{breaks}\SpecialCharTok{$}\NormalTok{brks}
\end{Highlighting}
\end{Shaded}

\begin{verbatim}
## [1]   4  22  45  68 122 253
\end{verbatim}

\textbf{Now that I have my breaks, I am going to use a facet map to
insets two maps into the same figure}

\begin{Shaded}
\begin{Highlighting}[]
\FunctionTok{print}\NormalTok{(}\StringTok{"Produce a facet map"}\NormalTok{)}
\end{Highlighting}
\end{Shaded}

\begin{verbatim}
## [1] "Produce a facet map"
\end{verbatim}

\begin{Shaded}
\begin{Highlighting}[]
\NormalTok{tm1 }\OtherTok{\textless{}{-}} \FunctionTok{tm\_shape}\NormalTok{(Accomodation\_contained) }\SpecialCharTok{+}
    \FunctionTok{tm\_polygons}\NormalTok{(}
      \AttributeTok{fill =} \FunctionTok{c}\NormalTok{(}\StringTok{"hotels\_n"}\NormalTok{, }\StringTok{"airbnbs\_n"}\NormalTok{),}
      \AttributeTok{fill.scale =} \FunctionTok{tm\_scale\_intervals}\NormalTok{(}\AttributeTok{values =} \StringTok{"brewer.blues"}\NormalTok{, }\AttributeTok{breaks=}\NormalTok{breaks}\SpecialCharTok{$}\NormalTok{brks),}
      \CommentTok{\# set the legend}
      \AttributeTok{fill.legend =} \FunctionTok{tm\_legend}\NormalTok{(}\AttributeTok{title=}\StringTok{"Accomodation count"}\NormalTok{,}
                              \AttributeTok{title.size=}\FloatTok{0.85}\NormalTok{,}
                              \AttributeTok{size=}\FloatTok{0.8}\NormalTok{,}
                              \CommentTok{\# plot outside of the main map}
                              \CommentTok{\#explained below}
                              \AttributeTok{position=}\FunctionTok{tm\_pos\_out}\NormalTok{(}\StringTok{"right"}\NormalTok{, }
                                                  \StringTok{"center"}\NormalTok{,}
                                                  \AttributeTok{pos.v =} \StringTok{"center"}\NormalTok{)),}
      \CommentTok{\# all facets share the same legend}
      \AttributeTok{fill.free =} \ConstantTok{FALSE}\NormalTok{)}\SpecialCharTok{+}
  \CommentTok{\# make 2 rows for the facets}
  \FunctionTok{tm\_facets}\NormalTok{(}\AttributeTok{nrow=}\DecValTok{2}\NormalTok{)}

\NormalTok{tm1}
\end{Highlighting}
\end{Shaded}

\begin{verbatim}
## Warning: Values have found that are less than the lowest break. They are
## assigned to the lowest interval
\end{verbatim}

\pandocbounded{\includegraphics[keepaspectratio]{Week_5_Practical_files/figure-latex/unnamed-chunk-17-1.pdf}}

\textbf{Now for the final map I want to \emph{alter the layout and add a
compass, scale bar, and north arrow}. Here is how to do that}

\begin{Shaded}
\begin{Highlighting}[]
\FunctionTok{print}\NormalTok{(}\StringTok{"produce final map"}\NormalTok{)}
\end{Highlighting}
\end{Shaded}

\begin{verbatim}
## [1] "produce final map"
\end{verbatim}

\begin{Shaded}
\begin{Highlighting}[]
\NormalTok{tm2 }\OtherTok{\textless{}{-}}\NormalTok{ tm1}\SpecialCharTok{+}
  \FunctionTok{tm\_layout}\NormalTok{(}\AttributeTok{panel.labels=}\FunctionTok{c}\NormalTok{(}\StringTok{"Hotels"}\NormalTok{, }\StringTok{"Airbnb"}\NormalTok{),}
            \AttributeTok{panel.show =} \ConstantTok{TRUE}
            \CommentTok{\#panel.label.bg.color = "transparent",}
\NormalTok{            )}\SpecialCharTok{+}
  
  \FunctionTok{tm\_compass}\NormalTok{(}\AttributeTok{type=} \StringTok{"arrow"}\NormalTok{,}
            \AttributeTok{size=}\FloatTok{1.8}\NormalTok{, }
            \AttributeTok{position =} \FunctionTok{tm\_pos\_out}\NormalTok{(}\StringTok{"left"}\NormalTok{, }
                                  \StringTok{"center"}\NormalTok{,}
                                  \AttributeTok{pos.h=} \SpecialCharTok{{-}}\FloatTok{0.05}\NormalTok{,}
                                  \AttributeTok{pos.v =}\FloatTok{0.72}\NormalTok{))}\SpecialCharTok{+}
  \FunctionTok{tm\_scalebar}\NormalTok{(}\AttributeTok{text.size =} \FloatTok{0.7}\NormalTok{,}
              \AttributeTok{width=}\DecValTok{10}\NormalTok{,}
              \AttributeTok{breaks=}\FunctionTok{c}\NormalTok{(}\DecValTok{0}\NormalTok{,}\DecValTok{10}\NormalTok{,}\DecValTok{20}\NormalTok{),}
              \AttributeTok{position =} \FunctionTok{tm\_pos\_out}\NormalTok{(}\StringTok{"left"}\NormalTok{,}
                                    \StringTok{"center"}\NormalTok{, }
                                    \AttributeTok{pos.h=}\FloatTok{0.075}\NormalTok{,}
                                    \AttributeTok{pos.v =} \FloatTok{0.68}\NormalTok{))}\SpecialCharTok{+}
  \CommentTok{\# we could use tm\_credits to place sub{-}titles like (A) or (B)}
  \CommentTok{\# on the map.}
  \FunctionTok{tm\_credits}\NormalTok{(}\StringTok{"(c) OpenStreetMap contrbutors and Air b n b"}\NormalTok{,}
             \AttributeTok{size=}\FloatTok{0.6}\NormalTok{,}
             \AttributeTok{position =} \FunctionTok{tm\_pos\_out}\NormalTok{(}\StringTok{"center"}\NormalTok{,}
                                   \StringTok{"bottom"}\NormalTok{,}
                                   \AttributeTok{pos.v =} \FloatTok{1.5}\NormalTok{,}
                                   \AttributeTok{pos.h=}\SpecialCharTok{{-}}\FloatTok{0.02}\NormalTok{))}

\NormalTok{tm2}
\end{Highlighting}
\end{Shaded}

\begin{verbatim}
## Warning: Values have found that are less than the lowest break. They are
## assigned to the lowest interval
\end{verbatim}

\begin{verbatim}
## For 'tm_scalebar()', 'breaks' and 'width' are not supposed to be used together; normally, setting the exact width is not needed when breaks have been specified.FALSE
\end{verbatim}

\pandocbounded{\includegraphics[keepaspectratio]{Week_5_Practical_files/figure-latex/unnamed-chunk-18-1.pdf}}

\textbf{Now I want to save the figure above}

\begin{Shaded}
\begin{Highlighting}[]
\FunctionTok{print}\NormalTok{(}\StringTok{"Save Map"}\NormalTok{)}
\end{Highlighting}
\end{Shaded}

\begin{verbatim}
## [1] "Save Map"
\end{verbatim}

\begin{Shaded}
\begin{Highlighting}[]
\FunctionTok{tmap\_save}\NormalTok{(tm2, }\StringTok{"Outputs/1\_facet.png"}\NormalTok{, }\AttributeTok{width =} \DecValTok{8}\NormalTok{, }\AttributeTok{height =} \DecValTok{6}\NormalTok{, }\AttributeTok{units=}\StringTok{"in"}\NormalTok{, }\AttributeTok{dpi =} \DecValTok{300}\NormalTok{)}
\end{Highlighting}
\end{Shaded}

\begin{verbatim}
## Warning: Values have found that are less than the lowest break. They are
## assigned to the lowest interval
\end{verbatim}

\begin{verbatim}
## For 'tm_scalebar()', 'breaks' and 'width' are not supposed to be used together; normally, setting the exact width is not needed when breaks have been specified.FALSE
\end{verbatim}

\begin{verbatim}
## Map saved to C:\Users\Owen Hughes\Desktop\CASA\UCL_Work\GIS_Term1\GIS_Git_Projects\Week_5\GIS\Week_5\Outputs\1_facet.png
\end{verbatim}

\begin{verbatim}
## Resolution: 2400 by 1800 pixels
\end{verbatim}

\begin{verbatim}
## Size: 8 by 6 inches (300 dpi)
\end{verbatim}

\textbf{Usually for this practical one would add the inset map of the UK
to show the study area, but that data download was not working so I will
skip this step this week.}

\end{document}
